\documentclass[11pt]{article}
\usepackage[margin=1in]{geometry}
\usepackage{amsmath,amssymb,amsthm}
\usepackage{xurl}
\usepackage{hyperref}
\sloppy
\let\oldus\_ \renewcommand{\_}{\oldus\allowbreak}
\newtheorem{theorem}{Theorem}
\newtheorem{lemma}[theorem]{Lemma}
\theoremstyle{definition}
\newtheorem{definition}[theorem]{Definition}

\title{Disproof of conjecture 00000004515:\\ the planar random Euclidean MST has no $(\log n)^{1/4}$ factor}
\author{Jackmeson1}
\date{2026-10-05}

\begin{document}
\maketitle

\section{The conjecture and the readings refuted}

\textbf{Statement (English, verbatim).} Definition: The random Euclidean MST is the minimum
spanning tree on iid points of the unit region. Conjecture: In dimension two the expected total
weight is a constant times sqrt(n) times the 1/4 power of log n, the exponent 1/4 is exact and
determined by the long-edge density formula. (log-power exponent 1/4 of the planar MST)
(The Chinese text states the same.)

\textbf{Objects.} Let $\mu$ be a probability law on $\mathbb{R}^2$ (the law of each point) and
let $X_1,\dots,X_n$ be iid with law $\mu$, i.e.\ $(X_1,\dots,X_n)$ has law $\mu^{\otimes n}$.
$W_n$ is the total Euclidean length of the minimum spanning tree of the complete graph on the
points, and $E_n = \mathbb{E}[W_n]$. Write $f(n)=\sqrt{n}\,(\log n)^{1/4}$ ($\log$ is the
natural logarithm; another base only changes the constant).

\textbf{Readings.} We refute each of the following, for \emph{every} law $\mu$ concentrated on
a bounded planar set (in particular the uniform law on the unit square or the unit disk):
\begin{itemize}
\item[(1)] $E_n \sim c\,f(n)$, i.e.\ $E_n/(c f(n)) \to 1$, for any real $c$;
\item[(2)] $E_n = c\,f(n)$ for all large $n$, for any $c>0$;
\item[(3)] $E_n = \Theta(f(n))$ (``the exponent $1/4$ is exact''): $c_1 f(n)\le E_n\le c_2 f(n)$
  eventually, with $c_1,c_2>0$;
\item[(4)] $E_n \ge c\,f(n)$ for infinitely many $n$, for any $c>0$.
\end{itemize}
Reading (4) is the weakest; (1) with $c>0$, (2) and (3) each imply it. The clause ``determined
by the long-edge density formula'' only explains where the exponent comes from; it falls with
the exponent. \emph{Not covered:} unbounded regions of unit area, power-weighted edge lengths
$|e|^\alpha$ with $\alpha\ne 1$, and $c=0$ in (2) (for a Dirac law $W_n=0$; the text's ``exponent
$1/4$ is exact'' excludes $c=0$).

\section{Definitions (as in the Lean file)}

\begin{definition}
For $p,q\in\mathbb{R}\times\mathbb{R}$, $d(p,q)=\sqrt{(p_1-q_1)^2+(p_2-q_2)^2}$ (\texttt{edist2};
Mathlib's \texttt{dist} on $\mathbb{R}\times\mathbb{R}$ is the sup metric, so the Euclidean formula
is written out). For $x:\mathrm{Fin}\,n\to\mathbb{R}^2$ and a simple graph $G$ on $\mathrm{Fin}\,n$,
$\mathrm{len}_x(G)=\sum_{\{i,j\}\in E(G)} d(x_i,x_j)$ (\texttt{graphLen}). The MST length is
\[ W(x)=\min_{T \text{ a tree on } \mathrm{Fin}\,n} \mathrm{len}_x(T) \]
(\texttt{mstLen}, an infimum over the finite set of trees; \texttt{mstLen\_attained} shows it is
attained for $n\ge1$). Every tree on the vertex set $\mathrm{Fin}\,n$ is a spanning tree of the
complete graph. The expectation is $E_n(\mu)=\int W\,d\mu^{\otimes n}$ (\texttt{expMST}, with
\texttt{Measure.pi}); \texttt{measurable\_mstLen} and \texttt{integrable\_mstLen} show that it is
a genuine finite expectation. The uniform law on the unit square is
$\mathrm{Leb}|_{[0,1]^2}$ (\texttt{uniformSq}, a probability measure).
\end{definition}

The four readings are the Lean predicates \texttt{AsympReading}, \texttt{EqReading},
\texttt{ThetaReading} and \texttt{LowerIOReading}, with $f$ = \texttt{prof}.

\section{The strip bound}

\begin{lemma}[\texttt{exists\_connected\_short}, \texttt{mstLen\_le\_unitSq}]
For $n\ge1$ points $x_1,\dots,x_n\in[0,1]^2$ there is a connected graph $G$ on them with
$\mathrm{len}_x(G)\le 3\sqrt n+2$. Hence $W(x)\le 3\sqrt n+2$.
\end{lemma}

\begin{proof}
Let $k=\lfloor\sqrt n\rfloor+1$, so $\sqrt n\le k\le\sqrt n+1$. The strip index of a height
$y\in[0,1]$ is $s(y)=\min(\lfloor ky\rfloor,k-1)$, so $s\le ky\le s+1$ and $0\le s\le k-1$.
Sort the points by the key $s(y)+x/2$ (\texttt{Tuple.sort}); since $x/2\in[0,1/2]$, this is the
lexicographic order (strip, then abscissa). Let $q_0,\dots,q_{n-1}$ be the sorted points and $G$
the path $q_0q_1\cdots q_{n-1}$ (the image of \texttt{pathGraph} under the sorting permutation, so
connected). For consecutive $p=q_j$, $q=q_{j+1}$ with strips $s\le t$:
if $s=t$ then $q_1\ge p_1$, so $|q_1-p_1|=q_1-p_1$; if $s<t$ then $|q_1-p_1|\le q_1-p_1+2\le
q_1-p_1+2(t-s)$. Also $|q_2-p_2|\le(t-s+1)/k$. Since $d\le|\Delta_1|+|\Delta_2|$,
\[ d(q_j,q_{j+1})\le F(j+1)-F(j),\qquad F(j)=(q_j)_1+(2+1/k)\,s(q_j)+j/k \]
(\texttt{step\_le}). Every edge of $G$ is such a pair, so telescoping gives
\[ \mathrm{len}(G)\le F(n-1)-F(0)\le 1+(2+1/k)(k-1)+(n-1)/k \le 3\sqrt n+2 \]
(\texttt{strip\_arith}: multiply by $k$ and use $k\le\sqrt n+1$, $n\le k\sqrt n$). A connected
graph contains a spanning tree $T\le G$ (Mathlib \texttt{Connected.exists\_isTree\_le}) and
$\mathrm{len}(T)\le\mathrm{len}(G)$ since lengths are nonnegative, so $W\le\mathrm{len}(G)$
(\texttt{mstLen\_le\_of\_connected}).
\end{proof}

\begin{lemma}[\texttt{mstLen\_le\_bounded}]
If $S\subset\mathbb{R}^2$ is bounded, there is $C\ge0$ with $W(x)\le C(3\sqrt n+2)$ for all
$n\ge1$ and all $x\in S^n$.
\end{lemma}

\begin{proof}
$S\subseteq[-R,R]^2$ for some $R>0$. The map $\phi(p)=((p_1+R)/2R,(p_2+R)/2R)$ sends $S$ into
$[0,1]^2$ and $d(p,q)=2R\,d(\phi p,\phi q)$ (\texttt{edist2\_sc}). Apply the previous lemma to
$\phi(x_i)$ and take $C=2R$.
\end{proof}

\section{The disproof}

\begin{theorem}[\texttt{ratio\_tendsto\_zero}, \texttt{not\_frequently\_lower}]
Let $\mu$ be a probability law such that $\mu$-almost every point lies in a bounded set $S$. Then $0\le E_n\le C(3\sqrt n+2)$ for $n\ge1$, and for every real $c$ and
every $\beta>0$,
\[ \frac{E_n}{c\,\sqrt n\,(\log n)^\beta}\longrightarrow 0 .\]
Consequently, for $c>0$, $E_n\ge c\sqrt n(\log n)^\beta$ holds for only finitely many $n$.
\end{theorem}

\begin{proof}
Almost surely all $X_i\in S$, so $0\le W\le C(3\sqrt n+2)$ a.s.\ and the bound integrates
(\texttt{expMST\_le}). For $c=0$ the quotient is $0$ by the convention $a/0=0$. For $c\ne0$ and
$n\ge3$, $|E_n/(c\sqrt n(\log n)^\beta)|\le 5C/(|c|(\log n)^\beta)\to0$ because
$(\log n)^\beta\to\infty$. If $E_n\ge c\sqrt n(\log n)^\beta$ for infinitely many $n$ with
$c>0$, the quotient would be $\ge1$ infinitely often, a contradiction.
\end{proof}

\begin{theorem}[\texttt{conjecture\_4515\_false}, \texttt{conjecture\_4515\_false\_uniform}]
For every such $\mu$, and in particular for the uniform law on $[0,1]^2$: (1) fails for every
real $c$; (2) fails for every $c>0$; (3) fails; (4) fails for every $c>0$.
\end{theorem}

\begin{proof}
(1): the quotient tends to $0\ne1$. (4): previous theorem with $\beta=1/4$. (2) and (3) each
give $E_n\ge c f(n)$ for all large $n$, hence for infinitely many $n$, contradicting (4).
\end{proof}

\textbf{Context.} The retrieved Wikipedia article ``Euclidean minimum spanning tree''
(\url{https://en.wikipedia.org/wiki/Euclidean_minimum_spanning_tree}) states: ``The total length
of the edges, for points in a unit square, is at most proportional to the square root of the
number of points.'' The proof above does not use it.

\section{Lean formalization}

File \texttt{lean/Conjecture4515/Basic.lean} (namespace \texttt{C4515}, only
\texttt{import Mathlib}, Mathlib v4.33.1). Main statements:
\begin{verbatim}
theorem conjecture_4515_false (mu : Measure (R x R)) [IsProbabilityMeasure mu]
    {S : Set (R x R)} (hS : Bornology.IsBounded S) (hmu : forall^ae p d mu, p in S) :
    (forall c : R, not AsympReading mu c) /\ (forall c > 0, not EqReading mu c) /\
      not ThetaReading mu /\ (forall c > 0, not LowerIOReading mu c)
theorem conjecture_4515_false_uniform :
    (forall c : R, not AsympReading uniformSq c) /\ (forall c > 0, not EqReading uniformSq c) /\
      not ThetaReading uniformSq /\ (forall c > 0, not LowerIOReading uniformSq c)
\end{verbatim}
(ASCII transcription of the Lean statements.) Supporting theorems: \texttt{mstLen\_le\_unitSq},
\texttt{mstLen\_le\_bounded}, \texttt{integrable\_mstLen}, \texttt{expMST\_le},
\texttt{ratio\_tendsto\_zero}, \texttt{not\_frequently\_lower}.

\textbf{Axiom audit.} \texttt{\#print axioms} for both main theorems reports only
\texttt{propext}, \texttt{Classical.choice}, \texttt{Quot.sound}. No \texttt{sorry}, no
\texttt{native\_decide}, no added axioms. Built with \texttt{lake build}.

\end{document}
